% !TEX root = quickstart.tex
\pagestyle{empty}

\begin{center}
    \includegraphics[width=119mm]{logo.eps}\\
    \vspace{4cm}
    \rule{\textwidth}{1.6pt}\vspace*{-\baselineskip}\vspace*{2pt}
    \rule{\textwidth}{0.4pt}\\[\baselineskip]
    {\fontsize{40}{50}\bfseries\scshape LibreS3 Quick Start Guide}
    \rule{\textwidth}{0.4pt}\vspace*{-\baselineskip}\vspace{3.2pt}
    \rule{\textwidth}{1.6pt}\\[\baselineskip]
\end{center}

%\definecolor{lightblue}{rgb}{.780,.886,.953}
\definecolor{lightblue}{rgb}{.835,.910,.965}
\lstset {
    escapeinside=\^\^,
    frame=single,
    breaklines=true,
    backgroundcolor=\color{lightblue},
}

\LARGE
\chapter{Introduction}

\indent LibreS3 is a robust Open Source implementation of the Amazon S3 service,
supporting a subset of the S3 REST API.

Standard S3 client libraries and tools can be used to access it (for example s3cmd, python-boto, DragonDisk, etc.) .

LibreS3 uses Skylable \SX as the storage backend, which provides data
deduplication and replication.\\

\section*{Useful links}
\begin{itemize}
    \item \url{http://lists.skylable.com}
    \item \url{https://bugzilla.skylable.com}
    \item \url{http://wiki.skylable.com}
    \item \url{http://www.skylable.com/products/libres3}
\end{itemize}

\chapter{Installation}

\section*{Requirements}

LibreS3 is regularly tested on Linux and FreeBSD. We recommend using the
binary packages from \url{http://www.skylable.com/download} if your platform
is supported. Otherwise, on most Unix platforms you can compile LibreS3
from source. You will need the following packages to be installed together with their
development versions:
\begin{itemize}
    \item OCaml (>= 3.12.1)
    \item camlp4 (matching your OCaml compiler version)
    \item OpenSSL
    \item PCRE C library
    \item GNU Make and m4
\end{itemize}
For example, on Debian run:
\small
\begin{lstlisting}
# apt-get install ocaml-native-compilers camlp4-extra libssl-dev libpcre3-dev make m4
\end{lstlisting}
\LARGE
On Fedora run:
\small
\begin{lstlisting}
# yum install ocaml /usr/bin/camlp4of /usr/bin/camlp4rf /usr/bin/camlp4 pcre-devel openssl-devel make m4 ncurses-devel
\end{lstlisting}
\LARGE

\newpage
\section*{Compilation}

Follow the standard installation procedure. In this guide we will install
LibreS3 into \verb+/opt/libres3+.
\small
\begin{lstlisting}
$ ./configure --prefix=/opt/libres3 && make
# make install
\end{lstlisting}
\LARGE

\chapter{Configuration}

\section*{Requirements}

LibreS3 by default listens on ports 8008 and 8443, which need to be available on a given
IP address.

LibreS3 connects to the SX cluster via HTTP(S). You can run LibreS3 and SX
on the same or different hosts.

\subsection*{DNS zone entry}

S3 buckets require a wildcard A record pointing to the IP address (1.2.3.4 below) of the
host running LibreS3, for example:
\small
\begin{lstlisting}
*.libres3.example.com. A 1.2.3.4
\end{lstlisting}
\LARGE

In case you don't have control over the DNS you'll have to modify the \verb|/etc/hosts| file
of each client machine and add a line for each bucket you want to access:
\small
\begin{lstlisting}
libres3.example.com 1.2.3.4
bucket1.libres3.example.com 1.2.3.4
bucket2.libres3.example.com 1.2.3.4
\end{lstlisting}
\LARGE

\section*{Setting up a LibreS3 node}

Setting up LibreS3 is as simple as running the interactive tool \verb+libres3_setup+.
If you provide the path to an existing \SX cluster configuration file created
by \verb+sxsetup+, most of the settings will be done automatically.

Please make sure the default volume replica count setting is less or equal
to the number of nodes in the SX cluster.

In the examples below we assume you
have an \SX cluster running on \verb|sx.example.com| and you want to run
LibreS3 on \verb|libres3.example.com|.

Example setup with sxsetup.conf:
\small
\begin{lstlisting}
# /opt/libres3/sbin/libres3_setup  --sxsetup-conf /opt/sx/etc/sxserver/sxsetup.conf
Successfully loaded SX configuration from '/opt/sx/etc/sxserver/sxsetup.conf'
Successfully obtained admin key
Successfully obtained sx server ip/dns name
Successfully obtained run as user
Successfully obtained run as group
Successfully obtained ssl key file
Successfully obtained ssl certificate file

S3 (DNS) name: libres3.skylable.com

Default volume replica count: 1

Generating '/opt/libres3/etc/libres3/libres3.conf'
Updating '/opt/libres3/etc/libres3/libres3-insecure.sample.s3cfg'
Successfully loaded SX configuration from '/opt/libres3/etc/libres3/libres3-insecure.sample.s3cfg'
Updating '/opt/libres3/etc/libres3/libres3.sample.s3cfg'
Successfully loaded SX configuration from '/opt/libres3/etc/libres3/libres3.sample.s3cfg'

Do you want to start LibreS3 now? [Y/n] Y
Starting LibreS3
LibreS3 started successfully
\end{lstlisting}
\LARGE

Example without sxsetup.conf:
\footnote{you can use 'sxserver status' on the SX node to find out the required information}
\small
\begin{lstlisting}
# /opt/libres3/sbin/libres3_setup

Admin key: 0DPiKuNIrrVmD8IUCuw1hQxNqZeKUL6IZ26Lro8FqwNmCJa9pH5XkgAA

SX server IP/DNS name: 127.0.1.1

Run as user: nobody

Run as group: nogroup

SSL key file: /opt/sx/etc/ssl/private/sxkey.pem

SSL certificate file: /opt/sx/etc/ssl/certs/sxcert.pem

S3 (DNS) name: libres3.skylable.com

Default volume replica count: 1

Generating '/opt/libres3/etc/libres3/libres3.conf'
File '/opt/libres3/etc/libres3/libres3.conf' already exists, overwriting
Updating '/opt/libres3/etc/libres3/libres3-insecure.sample.s3cfg'
Successfully loaded SX configuration from '/opt/libres3/etc/libres3/libres3-insecure.sample.s3cfg'
Updating '/opt/libres3/etc/libres3/libres3.sample.s3cfg'
Successfully loaded SX configuration from '/opt/libres3/etc/libres3/libres3.sample.s3cfg'

Do you want to start LibreS3 now? [Y/n] Y
Starting LibreS3
LibreS3 started successfully
\end{lstlisting}
\LARGE

To start/stop LibreS3\footnote{LibreS3 and SX will communicate using SSL by
default. For debugging purposes you can configure SX with 'sxsetup --no-ssl' and then you have to start LibreS3 with --no-ssl}:

\small
\begin{lstlisting}
# /opt/libres3/sbin/libres3 start
Starting LibreS3
LibreS3 started successfully
# /opt/libres3/sbin/libres3 status
--- LibreS3 STATUS ---
libres3 is running (PID 26780)

--- LibreS3 INFO ---
SSL private key: /opt/sx/etc/ssl/private/sxkey.pem
LibreS3 logfiles: /opt/libres3/var/log/libres3
# /opt/libres3/sbin/libres3 stop
Sending SIGTERM to 26780
Waiting for  26780
\end{lstlisting}
\LARGE

If the server doesn't start, please check the log files for details.

That's it - your LibreS3 cloud storage is already up and running!
You can now connect to it with your favourite S3 client.

\chapter{Client configuration}
\section*{s3cmd}

You can use the generated s3cfg config file
\footnote{\verb|/opt/libres3/etc/libres3/libres3.sample.s3cfg|}
 or configure s3cmd from scratch.
Below we assume that your LibreS3
is running on \verb|"libres3.example.com"| and it supports SSL.
The important s3cmd configuration settings are:

\small
\begin{lstlisting}
^\textbf{use\_https}^ True
^\textbf{host\_base}^ libres3.example.com:8443
^\textbf{host\_bucket}^ %(bucket)s.libres3.example.com:8443
^\textbf{access\_key}^ <your-sx-username>
^\textbf{secret\_key}^ <your-sx-key>
\end{lstlisting}
\LARGE

In case you don't use SSL, please use the port 8008 instead of 8443, and set
\verb|"use_https"| to \verb|"False"|. Once you've configured s3cmd check
that it properly connects to LibreS3:
\small
\begin{lstlisting}
$ s3cmd ls --debug 2>&1 | grep host
\end{lstlisting}
\LARGE

Supported s3cmd commands:
\begin{description}
    \item[Bucket] \verb|mb|, \verb|rb|, \verb|ls|, \verb|la|, \verb|du|
    \item[Object] \verb|put|, \verb|get|, \verb|del|, \verb|sync|, \verb|info|,
        \verb|cp|, \verb|mv|
\end{description}

\section*{Python-boto}

S3 clients using Python boto are configured in \verb|~/.boto|.
A typical configuration looks as follows:
\small
\begin{lstlisting}
[Credentials]
aws_access_key_id=<your-sx-username>
aws_secret_access_key=<your-sx-key>
s3_host=libres3.example.com:8443
s3_port=8443
[Boto]
is_secure = True
\end{lstlisting}
\LARGE

Note that setting \verb|"s3_host"| will override the hostname you give to
applications on the command-line. If you are using an application that allows
setting the S3 hostname on the command-line, you might want to use that instead.

\section*{Other clients}
For information on other clients please refer to our wiki:
\url{http://wiki.skylable.com/wiki/LibreS3_Clients}

\chapter{Troubleshooting}
In case you face a problem connecting to your LibreS3 server, please check
the log files located at:\\ \verb|/opt/libres3/var/log/libres3/*.log|\\
You can also enable logging of full HTTP requests with:
\small
\begin{lstlisting}
# export LIBRES3_DEBUG=1
# /opt/libres3/sbin/libres3 restart
\end{lstlisting}
\LARGE
For more information and FAQ please visit \url{http://wiki.skylable.com}.

If you can't find your solution there, please subscribe to our mailing list
at \url{http://lists.skylable.com} and post about your issues.

